% CLAUDE REPORT::B3-COLLAPSE-0001::2026-11-11T14:52-08:00
\documentclass[11pt]{article}
\usepackage{fontspec}
\usepackage{amsmath,amssymb,amsthm}
\usepackage[margin=1in]{geometry}
\usepackage{hyperref}
\newtheorem{theorem}{Theorem}
\newtheorem{proposition}[theorem]{Proposition}
\newtheorem{corollary}[theorem]{Corollary}
\theoremstyle{definition}
\newtheorem{definition}[theorem]{Definition}
\theoremstyle{remark}
\newtheorem{remark}[theorem]{Remark}
\newcommand{\Wbar}{\overline{W}}
\title{Which proof system makes the B3 bridge weakest?\\
\large CLAUDE B3-COLLAPSE-0001 (scheduled routine, worst-case derivation directive)}
\author{Claude (scheduled routine; not one of LSPV-1..5)}
\date{2026-11-11 14:52 PST}
\begin{document}
\maketitle

\noindent\textbf{Status summary.} P versus NP remains \textbf{OPEN}. This report proves no lower bound on SAT decision time. It classifies the bridge
\[
\mathrm{B3}(R):\qquad \mathrm{SAT}\in\mathrm{P}\ \Longrightarrow\ \exists C,k\ \forall N:\ \Wbar_R(N)\le C(N+1)^k,
\qquad
\Wbar_R(N)=\max_{\substack{F\in\mathrm{UNSAT}\\|\langle F\rangle|\le N}}\ \min_{\pi:F\vdash_R\bot}\operatorname{cost}_R(\pi),
\]
according to how strong the fixed system $R$ is. Findings:
\begin{enumerate}
\item (KNOWN THEOREM; already in corpus) For the multisign-hypergraph reduction system without auxiliary variables, the lower-bound half $\Wbar_R(N)\ge 2^{\Omega(N/\log N)}$ is already a theorem (Tseitin/expander, Urquhart 1987; Ben-Sasson--Wigderson 2001), so $\mathrm{B3}(R)$ is \emph{equivalent} to P$\ne$NP. The bridge is not weaker than the target.
\item (PROVED/EXACT, elementary; new to this corpus as an equivalence) Under SAT$\in$P, $\Wbar_R$ is polynomial \emph{iff} $R$ is p-optimal. Hence $\mathrm{B3}(R)\iff(\mathrm{P}=\mathrm{NP}\Rightarrow R \text{ is p-optimal})$.
\item (KNOWN THEOREM, first cited here) P$=$NP implies E$=$NE by padding, and E$=$NE implies that a p-optimal proof system exists (Krajíček--Pudlák 1989). So in the world where B3 has content, a system satisfying it exists; the question is whether the \emph{preselected} $R$ is one.
\item (KNOWN THEOREM, Cook 1975; already in corpus LSPV-1/357) For $R=$ Extended Frege (equivalently the hypergraph system \emph{with} auxiliary definitions, i.e.\ Extended Resolution), B3 holds for every polynomial-time SAT algorithm whose correctness is provable in PV. Full $\mathrm{B3}(\mathrm{EF})$ is OPEN, and so is the EF lower bound.
\end{enumerate}
Recommendation: B3 is monotone in $R$ (Proposition~\ref{p:mono}), so the weakest form of the bridge is obtained by taking $R$ as strong as possible. Work on B3 should use $R\supseteq$ EF. Systems with known lower bounds should be used only as controls.

\section{Literature preflight (this run)}
Obstacle 1: the bridge B3 for a fixed $R$. Obstacle 2: lower bounds for systems that have extension variables. Searched and checked:
\begin{itemize}
\item S.~Cook, \emph{Feasibly constructive proofs and the propositional calculus}, STOC 1975: a PV proof of $\forall x\,\varphi(x)$ gives polynomial-size EF proofs of the length-$n$ instances. (KNOWN THEOREM.)
\item S.~Cook, R.~Reckhow, J.~Symb.~Logic 44 (1979): NP$=$coNP iff some propositional proof system is polynomially bounded. (KNOWN THEOREM.)
\item J.~Krajíček, P.~Pudlák, J.~Symb.~Logic 54 (1989) 1063--1079: E$=$NE implies p-optimal proof systems exist; NE$=$coNE implies optimal ones exist. Köbler--Messner--Torán (Inf.~Comput.~2003) weaken the hypotheses to EE$=$NEE and NEE$=$coNEE. (KNOWN THEOREM.)
\item J.~Krajíček, survey statement: PV proves P$=$NP iff PV proves NP$=$coNP iff PV defines a polynomial-time algorithm that outputs a satisfying assignment or an EF refutation. A superpolynomial EF lower bound would make P$\ne$NP consistent with PV. (KNOWN THEOREM.)
\item P.~Pudlák, \emph{Reflection principles, propositional proof systems, and theories} (arXiv:2007.14835, 2020): reflection formulas and their relation to p-simulation. (KNOWN THEOREM / survey.)
\item Haken 1985 (PHP, Resolution), Urquhart 1987 (Tseitin, Resolution), Ben-Sasson--Wigderson 2001 (width), Ajtai 1988 / Beame--Impagliazzo--Pitassi 1993 / Krajíček--Pudlák--Woods 1995 (bounded-depth Frege), Pudlák 1997 (Cutting Planes), Razborov 1998 (Polynomial Calculus). (KNOWN THEOREMS.)
\item Topology handoff 2026-11-09 and its six cited works: read. They concern NP-hardness of link decision problems. Under Theorem~\ref{t:collapse} such results bear only on families, not on B3, so this run makes no topological claim.
\end{itemize}

\paragraph{Barrier-transfer classification.} Optimization/OR duality (an LP lower bound certifies optimality of a given solution and says nothing about runtime) has the same shape as ``minimum derivation cost'' versus ``decision time''. This is a \emph{useful structural analogy}. The missing transfer theorem is exactly B3. Post-quantum ledger item 5 (worst-case/average-case reduction architecture): this is a \emph{failed transfer} here. B3 is a worst-case statement about a fixed system and needs no distributional reduction. The other PQ items do not apply.

\section{Setting}
\begin{definition}
A \emph{refutation system} $R$ is a polynomial-time relation $V_R(\langle F\rangle,\pi)$ that is sound (if it accepts, then $F\in$ UNSAT) and complete. It comes with a cost function $\operatorname{cost}_R$ that is \emph{polynomially tied} to bit length: there is a polynomial $q$ such that $|\pi|\le q(\operatorname{cost}_R(\pi)+|\langle F\rangle|)$ and $\operatorname{cost}_R(\pi)\le q(|\pi|)$. $R$ is \emph{p-bounded} if $\Wbar_R\le\mathrm{poly}$. $R$ \emph{p-simulates} $Q$ if a polynomial-time $g$ maps every $Q$-refutation of $F$ to an $R$-refutation of $F$. $R$ is \emph{p-optimal} if it p-simulates every refutation system.
\end{definition}
Unit-cost-per-distinct-node DAG costs for the multisign hypergraph calculus are polynomially tied when node labels have polynomial bit size. Systems with exponentially compressed nodes violate the tie and must be handled separately; the Time-Space auditor's role covers this.

\begin{definition}[Multisign hypergraph reduction system $H$]
Vertices are variables. A hyperedge is a set of signed vertices (a clause). A move merges hyperedges $e_1\ni x^{+}$ and $e_2\ni x^{-}$ into $(e_1\cup e_2)\setminus\{x^{\pm}\}$ (resolution), optionally followed by enlarging the result (weakening). The minimal reduced logical state of an UNSAT instance is the empty hyperedge. $\operatorname{cost}_H$ counts distinct derived hyperedges in the DAG, and the best derivation is granted for free. $H^{\mathrm{ext}}$ also allows a fresh vertex $z$ together with the three hyperedges defining $z\leftrightarrow(a\wedge b)$.
\end{definition}

\section{Results}
\begin{proposition}[lower-bound half for $H$; KNOWN THEOREM]\label{p:lb}
$W_H(F)$ equals the minimum DAG-Resolution refutation size of $F$, since weakening never shortens Resolution. For odd-charge Tseitin 3-CNFs on 3-regular expanders with $n$ vertices, $N=\Theta(n\log n)$ and $W_H\ge 2^{\Omega(n)}$. Hence $\Wbar_H(N)\ge 2^{\Omega(N/\log N)}$ along an infinite explicit family of 3-CNFs. The same holds for every $R$ that Resolution, Cutting Planes, Polynomial Calculus or bounded-depth Frege p-simulates, with weaker exponents (PHP and Tseitin lower bounds for those systems).
\end{proposition}

\begin{theorem}[B3 collapse; PROVED/EXACT, elementary; already recorded in corpus as a meta-barrier]\label{t:collapse}
If $\Wbar_R$ is unconditionally superpolynomial, then $\mathrm{B3}(R)\iff\mathrm{SAT}\notin\mathrm{P}$.
\end{theorem}
\begin{proof}
If SAT$\notin$P, B3 holds vacuously. If B3 holds and SAT$\in$P, then $\Wbar_R$ is polynomial, which contradicts the hypothesis. So B3 implies SAT$\notin$P.
\end{proof}
\noindent Consequently, for $R=H$ every remaining difficulty lies in B3, and proving $\mathrm{B3}(H)$ \emph{is} proving P$\ne$NP. No ``free best sequence'' convention changes this, because the minimum inside $W_H$ already grants the best sequence.

\begin{proposition}[PROVED/EXACT]\label{p:opt}
Assume SAT$\in$P, and let $R$ be a refutation system with polynomially tied cost. The following are equivalent: (a) $\Wbar_R$ is polynomially bounded; (b) $R$ is p-bounded in bit length; (c) $R$ is p-optimal.
\end{proposition}
\begin{proof}
(a)$\Leftrightarrow$(b) follows from the cost tie.
(b)$\Rightarrow$(c): suppose every UNSAT $F$ with $|\langle F\rangle|=N$ has an $R$-refutation of length at most $p(N)$. The language $L=\{(\langle F\rangle,1^m,u):\exists v,\ |uv|\le m,\ V_R(\langle F\rangle,uv)\}$ is in NP, hence in P. Starting from the empty prefix, extend $u$ one bit at a time while staying in $L$, with $m=p(N)$, and test $V_R$ at each length. This finds a refutation in polynomial time without using any $Q$-proof. So $g(F,\rho)=$ (the found refutation) p-simulates every $Q$.
(c)$\Rightarrow$(b): let $A$ be a polynomial-time SAT decider and let $Q_A(\langle F\rangle,\rho)$ accept iff $\rho=\langle F\rangle$ and $A$ rejects $F$. $Q_A$ is sound and complete, and every refutation in it has length $N$. If $R$ p-simulates $Q_A$, the translation outputs $R$-refutations of length $\mathrm{poly}(N)$.
\end{proof}

\begin{corollary}
$\mathrm{B3}(R)\iff\bigl(\mathrm{SAT}\in\mathrm{P}\Rightarrow R\text{ is p-optimal}\bigr)$. With Krajíček--Pudlák 1989 and padding (P$=$NP$\Rightarrow$E$=$NE), SAT$\in$P implies that some p-optimal $R^\ast$ exists, so the consequent of B3 can be satisfied. The open question is whether a \emph{fixed, preselected} $R$ satisfies it. The KP89 system depends on the hypothesis and is not of this kind.
\end{corollary}

\begin{proposition}[monotonicity; PROVED/EXACT]\label{p:mono}
If $R'$ p-simulates $R$ (both with polynomially tied cost), then $\mathrm{B3}(R)\Rightarrow\mathrm{B3}(R')$. Conversely, a superpolynomial lower bound for $R'$ implies one for $R$. So B3 gets weaker and the lower bound gets harder as $R$ grows.
\end{proposition}
\begin{proof}
Translate refutations: $\Wbar_{R'}(N)\le\mathrm{poly}(N+\Wbar_R(N))$.
\end{proof}

\begin{remark}[conditional dichotomy; the core of this run]\label{r:dich}
No choice of $R$ makes both halves easier than the target. If $R$ is weak enough to have a known lower bound, B3 is equivalent to P$\ne$NP (Theorem~\ref{t:collapse}). If $R\supseteq$ EF, B3 has unconditional partial support, but the lower bound is a famous open problem: a superpolynomial EF lower bound alone would already show that PV does not prove P$=$NP. Moving the multisign hypergraph program from $H$ to $H^{\mathrm{ext}}$ (which is p-equivalent to Extended Resolution and hence to EF) is the only change that turns B3 from a reformulation into a candidate weaker step.
\end{remark}

\section{$\mathrm{B3}(\mathrm{EF})$ and the exact residual gap}
By Cook 1975, if $\mathrm{PV}\vdash\forall F,\alpha\,(A(F)=0\to\neg\mathrm{Sat}(F,\alpha))$, then the tautologies $\|\mathrm{RS}_A\|_n$ have polynomial-size EF proofs. Substituting $\langle F\rangle$ and evaluating $A(F)$ inside EF in $\mathrm{poly}(T_A)$ lines gives polynomial EF refutations of every UNSAT $F$. So $\mathrm{B3}(\mathrm{EF})$ holds for every PV-provably correct polynomial-time SAT algorithm. (KNOWN THEOREM.) The residual gap is the statement
\[
\mathrm{SAT}\in\mathrm{P}\ \Longrightarrow\ \exists A\in\mathrm{P}\text{ deciding SAT with }\ \|\mathrm{RS}_A\|_n\text{ having }\mathrm{poly}(n)\text{-size EF proofs},
\]
which is OPEN. LSPV-1/357's Gödelian firewall shows that a \emph{given} correct $A$ can be replaced by an extensionally equal $A'$ whose soundness is unprovable. That does not refute the residual statement, which quantifies existentially over $A$.

\section*{CORPUS REVIEWED}
Repository tree survey; \texttt{AGENTS.md}; root \texttt{main.tex} (section index and the R\_mix341/B3 comments near lines 4500--4800); \texttt{research/coordination/2026-11-02\_1204\_WORST\_CASE\_MINIMUM\_DERIVATION\_FIVE\_ROLE\_REFOCUS.md}; LSPV-1/335 worst-case core and generator barrier; LSPV-1/357 Gödel reflection B3 firewall and LSPV-1/361 ER reflection reports; LSPV-5 462--465 prewrite records (B3 $\iff$ $Q_A\le_W R$); the 2026-11-09 knot/link handoff; the latest handoffs dated 2026-11-11 (LSPV-1/581, LSPV-2/565, LSPV-3/578, LSPV-4/551); the Claude world ledger L01--L40.

\section*{DUPLICATION / CONTRADICTION AUDIT}
Theorem~\ref{t:collapse} duplicates the corpus meta-barrier ``B3(Res) iff P$\ne$NP because $\Wbar_{\mathrm{Res}}$ is superpolynomial (Haken)''. It is recorded here only to make the ledger complete and is not claimed as new. ``B3 iff $Q_A\le_W R$'' (LSPV-5) is the (c)$\Rightarrow$(b) half of Proposition~\ref{p:opt}. The corpus notes that p-optimality ``would suffice''; Proposition~\ref{p:opt} adds the converse (necessity under SAT$\in$P), and the corpus did not previously cite Krajíček--Pudlák E$=$NE. The Cook-1975 transfer agrees with LSPV-1/357. No contradiction was found.

\section*{WHAT CHANGED RELATIVE TO ROOT MAIN}
Nothing in root \texttt{main.tex} was edited. This routine owns no ownership region there, and LSPV-5 remains synthesis arbiter. New world-ledger rows L41--L46 were added (\texttt{research/agents/claude/world-ledger/}). Their automatic categories are L41 Theorem, L42 Equivalent to P$\ne$NP, L43--L45 Theorem, and L46 Half-win (true if P$\ne$NP). All six pass the consistency check.

\section*{ADOPTED / REJECTED / RETIRED PRIOR COMPONENTS}
Adopted: the refocus definitions of $\Wbar_R$, B1/B2/B3, and the cost conventions; the LSPV-1/335 MU-core and generator lemmas; and LSPV-1/357. Rejected as a route to a \emph{weaker} bridge: any fixed $R$ that a system with a known lower bound p-simulates. This is a reclassification of status, not a falsification of mathematics.

\section*{STALE OR OVERSTATED CLAIMS FOUND}
Phrasing in the directive such as ``prioritize the weakest sufficient bridge'' together with ``establish a superpolynomial lower bound in a fixed sufficiently powerful system'' can be read as treating the two halves as independent. By Theorem~\ref{t:collapse} and Remark~\ref{r:dich} they are not independent: for weak $R$ the bridge is the entire problem. This is a framing caution, not an error in any verified report.

\section*{Relativization / model audit}
Theorem~\ref{t:collapse} and Propositions~\ref{p:opt} and~\ref{p:mono} relativize. They are statements about proof systems and say nothing about machine models. No resource model, wrapper or SAME\_REALIZATION claim is made.

\noindent P versus NP remains OPEN; every load-bearing arrow (a lower bound for an $R\supseteq$ EF, and $\mathrm{B3}(R)$ for the same $R$) is OPEN.

\end{document}
% END_OF_VERIFIED_REPORT::CLAUDE::B3-COLLAPSE-0001
